\section{Historical Notes on the Harmonic Sequence}

One might wonder why the sequence
\[
 1, \frac 1 2, \frac 1 3, \frac 1 4, \frac 1 5, \dots
\]
is called \emph{harmonic}. Does it have anything to do with musical harmony?

According to tradition, Pythagoras once visited a blacksmith's workshop, where
he heard the sounds of hammers striking anvils. He noticed that some hammers
produced sounds that were harmonious with each other, while others did not.
Intrigued, he studied the phenomenon by attempting to produce sounds of
different frequencies using strings with varying tensions and lengths. He
discovered that sounds perceived as harmonious corresponded to
\emph{commensurable} frequencies, meaning they could be expressed as ratios of
integers.

Naturally, the sensation of "harmonicity" is subjective and produced by our
brain, but we can hypothesize why a frequency that is a multiple of another is
perceived as harmonic. Consider a string instrument: plucking the string induces
a vibration that can be described as a superposition of sinusoidal waves (see
the theory of Fourier series in chapter~\ref{sec:convergenza_integrale}). Since
the string is fixed at both ends, one possible oscillation is where the string
assumes the shape of a sinusoid with a half-period equal to the length of the
string itself. However, other superimposed oscillations are also possible: for
example, an oscillation with a period equal to the length of the string,
creating a node (fixed point) at the center. If we touch the string at one-third
of its length, we force the creation of a node at that point, preventing the
formation of the two fundamental waves and allowing only the \emph{harmonics}
corresponding to oscillations with at least one node at that point to be heard.
Essentially, the sound produced is a superposition of pure sinusoidal waves,
with frequencies that are integer multiples of the fundamental frequency. As the
frequencies increase, their intensity rapidly decreases; otherwise, the sum of
the sinusoids (the series) would diverge. The different intensity of the
harmonics determines the timbre of the musical instrument. For example, in a
flute, at the open end of the tube, there is an \emph{antinode} instead of a
node, and only the odd multiples of the fundamental frequency are obtained.
Since it is very difficult to obtain a pure sound, the presence of harmonics is
ubiquitous in natural sounds. It is therefore natural that two sounds with
frequencies that are multiples of each other are perceived as harmonious, while
sounds with frequencies in ratios far from small integers are perceived as
dissonant.

Based on these considerations, we can understand how musical notes are defined
in the Western tempered system. We choose as a reference the sound produced by a
string of a certain length, called the \emph{tonic}. If we halve the length of
the string, we obtain a sound with a fundamental frequency that is double: this
note is an \emph{octave} above the tonic (the reason for the name "octave" will
be clarified later). Our perception tells us that this sound is higher than the
previous one, but so similar that we tend to give it the same name. The same
happens if we double the length of the string (halving the frequency): we obtain
the same note, but lower, that is, an octave below the tonic.

Our perception of frequencies is therefore \emph{logarithmic} and not linear:
the transition from one frequency to double that frequency is perceived as the
same interval, regardless of the starting frequency. If we triple the frequency,
we obtain a \emph{different} note but very consonant with the tonic, called the
\emph{dominant}. The interval between the tonic and the dominant is called a
\emph{fifth} (we will see later the reason for the name). Since doubling or
halving the frequency does not change the perception of the note, the dominant
can be represented by the ratio $\frac 3 2$, intermediate between 1 and 2. The
harmonic relationship between tonic and dominant can be inverted by dividing the
frequency by three: bringing the frequency $\frac 1 3$ into the interval between
$1$ and $2$ (by multiplying by four, that is, ascending by two octaves), we
obtain the ratio $\frac 4 3$, called the \emph{subdominant} because it is lower
than $\frac 3 2$.

If we call the starting note Do (tonic), the notes with double or half frequency
will always be called Do, but on the higher and lower octaves relative to the
starting Do. The triple and one-third frequencies will be called Sol and Fa,
respectively. In the interval between two consecutive Do notes with frequencies
$f$ and $2f$, we will find a Fa with frequency $\frac 4 3 f$ and a Sol with
frequency $\frac 3 2 f$.

To better understand this structure, it is useful to switch to the logarithmic
scale: the logarithms of the ratios between frequencies become differences of
the logarithms. Calculating the base $2$ logarithms of the ratios with the
frequency of the reference Do note, the Do notes correspond to integers: the
reference Do to $0$, the octave above (double frequency) to $1$, the second
octave above corresponds to $2$, the lower octave to $-1$, and so on. The triple
frequency of the Sol notes, brought into the interval $[0,1]$, corresponds to
$\log_2\frac 3 2 \approx 0.585$, while that one-third of the Fa notes
corresponds to $\log_2 \frac 4 3 \approx 0.415$.

\begin{figure}
    \begin{center}
        \setlength{\tabcolsep}{1mm}% reduces the space between columns to 1mm
        \begin{tabular}{@{}l|ccccccccccccc@{}}
                        notes   &   Do   &   Reb           &   Re     &    Mib
            &   Mi          &    Fa    &  Solb/Fa\#                       &
            Sol    &  Lab           &   La          &   Sib         &   Si
            \\
                        fifths &   $0$  &  $-5$           &  $+2$    &    $-3$
            &  $+4$         &   $-1$   &  $-6/+6$                         &
            $+1$    &  $-4$          &  $+3$         &  $-2$         &  $+5$
            \\
                        freq   &   $1$  &$\frac{256}{243}$&$\frac
            98$&$\frac{32}{27}$&$\frac{81}{64}$&$\frac
            43$&$\frac{1024}{729}/\frac{729}{512}$&$\frac
            32$&$\frac{128}{81}$&$\frac{27}{16}$&$\frac{16}9$
            &$\frac{243}{128}$\\
                $\log_2$   &\small 0&\small .075 &\small  .170& \small  .245&
        \small  .340 & \small  .415& \small  .490 / .510         &\small
        .585&\small  .660&\small  .755  &\small  .830  &\small  .925    \\
                            temp &\small 0&\small .083 &\small  .167& \small
              .250& \small  .333 & \small  .417& \small  .500
              &\small  .583&\small  .667&\small  .750  &\small  .833  &\small
              .917    \\
        \end{tabular}
    \end{center}
    \caption{Table of notes, frequency ratios, and base 2 logarithms. The last row shows the equal (tempered) division of the octave.}
\end{figure}

We can now complete the scale of musical notes by adding the triple and
one-third frequencies of the Fa and Sol notes. Since the interval between Do and
Sol and between Fa and Do is called a fifth interval, the construction we are
describing is called the circle of fifths. The triple frequency of a note, on
the logarithmic scale (modulo octaves), corresponds to adding $0.585$ or
subtracting $0.415$. Ascending (by a fifth) from a Sol, we obtain a note with
frequency $\frac 9 8$ relative to the reference Do, whose logarithm has a
fractional part $\approx 0.170$; we call this note $Re$. Continuing in this way,
multiplying and dividing the frequencies by $3$, we obtain notes that are
distributed in the frequency interval between $f$ and $2f$. The frequencies that
are multiples of $3$ starting from Do give, in order: Sol, Re, La, Mi, Si, and
Fa\#. Dividing by 3, we obtain: Fa, Sib, Mib, Lab, Reb, and Solb. The
frequencies of the Fa\# and Solb notes are very close: $\frac{3^6}{2^9} =
\frac{729}{512} \approx 1.42$ and $\frac{2^{10}}{3^6} = \frac{1024}{729} \approx
1.40$. The base two logarithms of these fractions are $\log_2 \frac{729}{512}
\approx 0.510$ and $\log_2 \frac{1024}{729} \approx 0.490$. The two frequencies
are practically indistinguishable, and we can therefore identify the two notes
as the same. The difference between these frequencies is called the
\emph{musical comma}. Thus, we obtain 12 different notes, distributed almost
uniformly in the interval between the base frequency $f$ and its octave $2f$.
The correction of this small error can be distributed among the different notes
in various ways, giving rise to different tempering systems. The
\emph{well-tempered} scale (which became prevalent after Bach demonstrated its
versatility in composition) distributes the comma uniformly among the notes,
obtaining a scale in which the frequency ratio between two consecutive notes
(semitone) is exactly the twelfth root of two: $\sqrt[12]{2}\approx 1.06$. In
this way, the logarithms of the frequencies divide the interval $[0,1]$ into
twelve equal parts of width $\frac 1 {12} \approx 0.083$.

We have thus explained why 12 different musical notes were chosen: ultimately,
12 is the smallest number for which a power of 3 closely approximates a power of
2, given that $3^{12}$ is very close to $2^{19}$. The sequence of the smallest
powers of $3$ that best approximate the powers of $2$ (ordered by relative
error) is: 1, 2, 5, 12, 41, 53, 306, \dots. Thus, the choice of 12 notes is very
natural (even the pentatonic scale, with 5 notes, is common in many musical
cultures).

We have thus seen how the ratio $\frac 1 2$ determines the naming of the notes
(notes with double frequency are the same note) and how the ratio $\frac 1 3$
determines the choice of the musical scale and the number of distinct notes. The
subsequent harmonics are related to \emph{musical harmony}, that is, the
composition of chords and scales. The ratio $\frac 1 4$ does not generate new
notes, as like $\frac 1 2$, it represents the same musical note. The next ratio
is $\frac 1 5$, which approximately corresponds to an Mi (if, as above, we have
chosen Do as the reference note). The frequencies $1,3,5$ generate the major
chord, formed by the notes Do, Sol, Mi. Composing the major chords on the tonic
(Do, Mi, Sol), on the dominant (Sol, Si, Re), and on the subdominant (Fa, La,
Do), the seven notes of the major scale are obtained: Do, Re, Mi, Fa, Sol, La,
Si. For this reason, the interval between Do and Sol is called a \emph{fifth},
as it includes five of the seven notes, and the interval between Do and Fa is
called a \emph{fourth}. The interval that includes the entire scale is called an
\emph{octave}, as it includes the upper Do as the eighth note.